\section{Collision counts and microscopic windows}\label{sec:micro}
We prove the needed source extensions explicitly. In this section a
constant may depend on a requested fixed remainder accuracy.

\begin{lemma}[A passive incidence record]\label{lem:record}
The exact component expansion and unweighted variation estimates in
\Cref{prop:imports} also hold when a designated particle records its
number of accepted C contacts in an interval $I$.
\end{lemma}
\begin{proof}
In slab $i$ an accepted component has at most
$\Lambda_i-1+\Gamma$ edges. Thus each particle's total incidence record
is at most $B_\eps=\sum_i(\Lambda_i-1+\Gamma)$, a fixed power of
$|\log\eps|$. Normalize it by $B_\eps$.
In the source's autonomous partition identity, compatible autonomous
components reproduce the full accepted root history, hence this statistic.
The one-slab root-output kernel can therefore append it as one passive
coordinate. Spectator cancellation and the fresh-label identity are
unchanged. Pushforward has variation norm at most one; the crude operator
bound and stopped forest domination do not require derivatives or a
particular choice of output statistic.
Required record labels still continue to zero, while unmarked one-root
blocks are the same inputs $g_i$. Thus the source's finite regrouping and
coarse estimates apply with one bounded coordinate per required label,
not one coordinate per mesh point.
At total degree $k$, restoring the raw record costs $B_\eps^k$ on errors.
This is polylogarithmic, so exponentially small errors remain smaller than
every power, and arbitrary-power tails can be chosen after $k$.
On a retained diagram, raw mark powers cost only a fixed power of its
contact count, hence a fixed polynomial in size at complexity $K$.
\end{proof}

\begin{lemma}[Calendar localization]\label{lem:calendar}
For fixed factorial order, complexity and total contact-mark degree, the
sum of connected retained absolute marked integrals is at most
$C\mu^{1-h}|I|$, provided at least one marked contact is required to lie
in $I$. Nonphysical initial tails may instead be assigned arbitrary-power
errors.
\end{lemma}
\begin{proof}
First integrate out every common spatial translation and replace the
birth factors by their uniform envelope $C^m e^{-bE}$.
The remaining quotient coordinates are $(V,\Omega,t,y)$; there is no
remaining integration in an unrestricted translation variable.
Collapse identity interface gluings of continued lines.
For the positive comparison family, freely extend each global line
before its first and after its last prescribed contact. Redundant-contact
segments end at neighboring prescribed contacts or at infinity, not at
fixed physical lifetime interfaces. Keep the full chronological slot
order and the requirement that each redundant incoming intersection
exists; otherwise its branch weight is zero.

For $u\geq0$ set
\[
 t_e\longmapsto t_e+u,\qquad
 y_{pq}\longmapsto y_{pq}-u(v_p^{\rm birth}-v_q^{\rm birth}),
\]
leaving birth velocities and free normals fixed.
This quotient-coordinate map has determinant one. Reconstructing virtual
paths inductively shows that the new path at time $t$ is the old path at
$t-u$, up to the already integrated representative translation.
Merger normals, velocities and fluxes are unchanged. A redundant
incoming intersection shifts by $u$ with its normal and update unchanged.
The first shifted contact is later than the original one, so its fixed
birth anchor remains on its incoming free segment.
Endings and nominal slab interfaces are only constraints in the original
domain; virtual comparison paths may extend beyond them.
All global coordinates remain in the energy accounting, and every update
preserves total birth energy.

Put all absolute birth, ending, interface and nominal-slab restrictions,
including those on determined redundant times, into a domain $D$.
The Gaussian, fluxes and virtual chronological branch indicators form
a calendar-invariant core $W$.
If $t_*$ is a selected contact, $t_*(\Phi_u\theta)=t_*(\theta)+u$.
Changing variables, enlarging all shifted domains to $D_\ell^+$, and
averaging over $0\leq u\leq\ell$ gives
\begin{equation}\label{eq:calendar-average}
 \int_D\ind_{\{t_*\in I\}}W\dd\theta
 \leq\frac{|I|}{\ell}\int_{D_\ell^+}W^+\dd\theta ,
\end{equation}
because $\int_0^\ell\ind_{\{t_*\in I+u\}}\dd u\leq|I|$.
The changed interval $I+u$ is essential.

If $p=0$, take $\ell=\tau$. Nominal merger-time intervals have length at
most $2\tau$, and remain ordered within each slab.
Their volumes are $(2\tau)^{a_i}/a_i!$.
Retain the original global slot order until the final volume upper bound.
The same pair-speed, Gaussian-energy and factorial estimates apply to
these virtual histories: the subset of $n_i$ nominal labels has
$\sum_{p<q}|v_p-v_q|\leq n_i^{3/2}\sqrt E$.
Thus their size majorant is
$C_K\mu^{1-h}(1+n)^{d_K}(2C_b\tau\sqrt L)^n$.
Taking $L$ sufficiently large once absorbs the enlargement.

If $p\geq1$, take $\ell=\eps<\tau$. An initial-link ball enlarges to
radius $\eps(R_{K,D}+\sqrt{2E})$. The $p$ link volumes are at most
$C_{K,D}\eps^{3p}(1+E)^{3p/2}$.
The normalization remains $\mu^{1-h}\eps^p$.
Averaging in \eqref{eq:calendar-average} pays $\eps^{-1}$; the link gain
leaves $\eps^{p-1}\leq1$.
Absorb the fixed energy polynomial into half the birth Gaussian:
$(1+E)^{3p/2}e^{-bE}\leq C_{b,p}e^{-bE/2}$.
The fraction reserved is fixed independently of $p$, so the exponential
size base stays independent of complexity and moment order.
Nonphysical initial tails are erased with their inherited power cost
and chosen afterward.
Selecting a contact and restoring raw mark powers add only polynomial
size factors. Summing the geometric size tail proves the lemma.
\end{proof}

\begin{proposition}[Actual incidence cumulants]\label{prop:counts}
Let $A_I$ be the total number of collision incidences, two per accepted
C contact. For every fixed integers $k\geq1$ and fixed accuracy $M>0$,
\[
 |\kap_k(A_I)|\leq C_{k,M}(\mu|I|+\eps^M)
\]
uniformly in $I\subset[0,T]$.
\end{proposition}
\begin{proof}
Apply \Cref{lem:record}. In a connected retained term of total mark degree
$k$, a nonzero mark product selects at least one C contact in $I$.
\Cref{lem:calendar} bounds its factorial-order-$h$ contribution by
$C_{h,K,k}\mu^{1-h}|I|$. The ordinary-cumulant partition identity then
multiplies by $\mu^h$, giving $C\mu|I|$.
Choose complexity, initial tails and stopped-error powers after $k,M$,
allowing for the fixed factors $\mu^h$ and $B_\eps^k$.
The actual hierarchy-to-cumulant remainder is handled exactly as in
\Cref{lem:remainder}. The leading constant may depend on $M$ through $K$;
there is no assertion of a constant uniform over all accuracy choices.
\end{proof}

\begin{proposition}[Microscopic modulus]\label{prop:micro}
For every fixed $J\geq2$, $\lambda>0$ and compact smooth $\phi$,
\[
 \Pp\left\{\sup_{|t-s|\leq\eps^J}
       |\zeta_t^{\eps,p}(\phi)-\zeta_s^{\eps,p}(\phi)|>\lambda\right\}
       \leq C_{\phi,\lambda,J}\eps^2
\]
for sufficiently small $\eps$.
\end{proposition}
\begin{proof}
Put $\delta=\eps^J$ and cover the interval by $O(T/\delta)$ overlapping
windows of length $2\delta$. The raw fourth-moment identity and
\Cref{prop:counts}, with the accuracy fixed sufficiently large after $J$,
give
\[
 \E A_I^4\leq C\sum_{r=1}^4(\mu|I|)^r+O(\eps^M).
\]
For example accuracy $M+6$ in each cumulant absorbs the other at most
three factors of order $\mu$. Markov and the grid union bound yield
\[
 \Pp\{\max_I\eps A_I>\lambda\}
 \leq C_\lambda(\eps^2+\delta+\eps^{-2}\delta^2+
                         \eps^{-4}\delta^3)
       +O(\eps^{M+4-J})=O(\eps^2).
\]
Pathwise, with $D_\phi=\|v\cdot\nabla_x\phi\|_\infty$,
\[
 \osc_I\left(\eps\sum_i\phi(z_i^p(t))\right)
       \leq\eps D_\phi N|I|+2\|\phi\|_\infty\eps A_I.
\]
There are no state jumps when the pasted prehistory restarts.
The single global $N$ term needs no grid union:
$\E N^2\leq C\mu^2$ bounds its probability by
$C\eps^2\mu^2\delta^2=C\eps^{2J-2}\leq C\eps^2$.
The exact mean increment is deterministically at most
$C_\phi\eps^{-1}\delta+O(\eps^M)$ by the first incidence cumulant.
It is smaller than a fixed threshold for small $\eps$.
This proves the statement, with constants allowed to depend on $J$.
\end{proof}
